%% sections/slide05-headline-results.tex — SLIDE 5 — Headline results
%% =============================================================================
%% CONTENT BLUEPRINT — Headline (single most-important) result
%% =============================================================================
%% Time budget: 90 seconds — the largest single block in the talk.
%% This slide carries the verdict on Objective O1 (highlighted on slide 3).
%%
%% WHAT TO PUT ON THIS SLIDE:
%%
%%   ONE high-quality plot, large filling roughly 70 % of the slide area.
%%   Around it, a minimal annotation layer:
%%
%%     • Title strip at top: "Result 1 — [what was determined, with number]."
%%       Examples: "We measured n₂ of SF59 = 5.4×10⁻¹⁹ m²/W at 632.8 nm."
%%                 "We computed E_g(MoS₂) = 1.82 eV at 0 % strain (lit. 1.80)."
%%
%%     • Above the plot: a single sentence stating the headline number with
%%       units and uncertainty. The number must be the FIRST thing the
%%       examiner reads.
%%
%%     • Below the plot: 2–3 bullets explaining what they should read off
%%       the plot. ANNOTATE the key point on the figure using TikZ
%%       (\node, \draw arrow into a peak/valley) rather than relying on
%%       prose alone.
%%
%%     • At the very bottom: an orange-accented "Result" bar with a one-line
%%       summary. Use \resultbox{...}{...} from packages.tex.
%%
%% DO NOT PUT:
%%   • Two figures. If you have two equally-strong results, put them on
%%     separate slides: split slide 5's content into 5a and 5b — but be
%%     honest about the trade-off (less time per slide).
%%   • A long caption. Recreate the reading hierarchy: title $\to$ number
%%     $\to$ plot $\to$ explanation.
%%   • Error bars everywhere — pick the one most important source of
%%     uncertainty and annotate it on the figure.
%%
%% FIGURE SOURCING:
%%   • Pull from figures/headline-result.pdf (or .png) — a vector PDF is
%%     preferable so text remains crisp at projector resolution.
%%   • If you generated the figure in pgfplots directly in your thesis,
%%     you can lift the source — but simplify: drop minor-grid styling,
%%     drop per-point markers unless you have <20 points, axis ranges tight
%%     to the data.
%%
%% USE THE ANNOTATED-FIGURE PATTERN BELOW AS A STARTING POINT.
%% =============================================================================

\begin{frame}{Headline Result~\textemdash~[the single most important finding]}
  \begin{resultbox}{Result 1.%
    \timenote{Phrase as a sentence ending in a number with units.}}{%
    \itshape ``We measured / computed \highlight{[target]}\,=\,
    \highlight{[value]} $\pm$\,\highlight{[uncertainty]}\,[\si{unit}] under
    [conditions], comparing to literature value
    \highlight{[lit.value]} \highlight{[ref]}.''}
  \end{resultbox}

  \vspace{0.15cm}
  \begin{columns}[T, totalwidth=\textwidth]
    \begin{column}{0.58\textwidth}
      \centering
      \vspace{-2pt}
      % --- Example headline plot (REPLACE with your own figure) ---
      \begin{tikzpicture}
        \begin{axis}[
            width=\textwidth, height=5.6cm,
            xlabel={Parameter $x$ (units)},
            ylabel={Observable $Y$ (units)},
            xmin=0, xmax=10, ymin=0, ymax=2.5,
            grid=both,
            grid style={line width=.1pt, color=gray!10},
            major grid style={line width=.2pt, color=gray!25},
            legend pos=south east, legend cell align={left},
            legend style={font=\scriptsize, draw=sustGray!40}
          ]
          % Simulated curve
          \addplot[color=sustBlue, very thick, domain=0:10, samples=80]
              {1 + 0.6*sin(deg(x))};
          \addlegendentry{This work (simulation)}

          % Mock data points
          \addplot[only marks, mark=*, color=sustAccent, mark size=1.8pt]
              coordinates {(1.0, 1.55) (2.0, 1.78) (3.0, 1.62)
                          (4.0, 1.10) (5.0, 0.95) (6.0, 1.05)
                          (7.0, 1.42) (8.0, 1.71)};
          \addlegendentry{Experimental}

          % Annotation — head of the examiner
          \node[circle, fill=sustAccent, inner sep=1.5pt, label=above:
                {\scriptsize\bfseries\color{sustAccent}peak here}] at (axis cs:2, 1.78) {};
          \draw[sustAccent, -Latex, thick] (axis cs:2,1.78) -- (axis cs:6.6, 2.2);
        \end{axis}
      \end{tikzpicture}
    \end{column}

    \begin{column}{0.40\textwidth}
      \vspace{-4pt}
      \textbf{\color{sustBlueDark}Read this off the plot}
      \begin{itemize}\setlength\itemsep{2pt}
        \item Trend is monotone / peaked / oscillatory --- say which.
        \item Quantitative value at annotated marked feature.
        \item Comparison vs.\ literature.
        \item State ONE source of uncertainty.
      \end{itemize}
      \vspace{0.1cm}
      \timenote{Remember: this slide decides whether the examiner is sold.
      Answer questions explicitly: yes / no / part; with what percentage.}
    \end{column}
  \end{columns}
\end{frame}
